\documentclass{article}

% Packages
\usepackage{amssymb,amsthm,amsmath}
\usepackage{tikz-cd}
\usepackage{mathtools}
\usepackage{microtype}
\usepackage{needspace}
\usepackage{hyperref}
\usepackage{cleveref}

% Theorem environments
\theoremstyle{definition}
\newtheorem{example}{Example}[section]
\newtheorem{nonexample}{Non-Example}[section]
\newtheorem{definition}{Definition}[section]

\theoremstyle{remark}
\newtheorem{remark}{Remark}[section]

\theoremstyle{plain}
\newtheorem{proposition}{Proposition}[section]
\newtheorem{theorem}{Theorem}[section]
\newtheorem{corollary}{Corollary}[section]
\newtheorem{exercise}{Exercise}[section]

% Mathematical notation
\DeclareMathOperator{\Spec}{Spec}

\newcommand{\N}{\mathbb{N}}
\newcommand{\Z}{\mathbb{Z}}
\newcommand{\Q}{\mathbb{Q}}
\newcommand{\R}{\mathbb{R}}
\newcommand{\C}{\mathbb{C}}
\newcommand{\A}{\mathbb{A}}
\newcommand{\F}{\mathbb{F}}
\renewcommand{\P}{\mathbb{P}}

\begin{document}

% Number the environments as 14.1, 14.2, ...
\setcounter{section}{13}
\stepcounter{section}

\section*{Lecture 14: Morphisms to Projective Space}

\subsection*{Locally split injections and fibres}

We have been trying to generalise the description of
\(\operatorname{Mor}(\Spec k,\P^n_\Z)\) to a description of
\(\operatorname{Mor}(X,\P^n_\Z)\). Recall that the proposed correspondence is
\[
    \operatorname{Mor}(X,\P^n_\Z)
    \longleftrightarrow
    \{(\mathcal{L},\iota)\}/\cong,
\]
where \(\mathcal{L}\) is a locally free \(\mathcal{O}_X\)-module of rank 1
and \(\iota:\mathcal{L}\to\mathcal{O}_X^{\oplus(n+1)}\) is a locally split
injection. Here an isomorphism of pairs must commute with the maps into
\(\mathcal{O}_X^{\oplus(n+1)}\).

For \(\Spec k\), a point of projective space is obtained by choosing a line
in \(k^{n+1}\). For an arbitrary scheme \(X\), we would similarly like to
choose a line at every point of \(X\), with these lines varying locally in
a coherent way. However, an injection of sheaves need not give an injection
after evaluation at every point. This is why we require a locally split
injection, rather than just an injection of sheaves.

Let us examine what this condition means for a map
\[
    \iota:\mathcal{L}\longrightarrow\mathcal{O}_X^{\oplus(n+1)}.
\]
Since \(\mathcal{L}\) is locally free of rank 1, we can choose an open set
\(U\subseteq X\) and a trivialisation
\(\tau:\mathcal{O}_U\xrightarrow{\sim}\mathcal{L}|_U\).
Restriction commutes with finite direct sums, so
\[
    \bigl(\mathcal{O}_X^{\oplus(n+1)}\bigr)|_U
    \cong\mathcal{O}_U^{\oplus(n+1)}.
\]
Thus, in this trivialisation, \(\iota|_U\) becomes the map
\[
    \phi=(\iota|_U)\circ\tau:
    \mathcal{O}_U\longrightarrow\mathcal{O}_U^{\oplus(n+1)}.
\]
The condition that \(\iota\) is locally split means that we can choose an
open cover by such sets \(U\) on each of which \(\phi\) is a split injection.

\begin{proposition}\label{prop:lsi-section-criterion}
    Let \(\phi:\mathcal{O}_U\to\mathcal{O}_U^{\oplus(n+1)}\) be an
    \(\mathcal{O}_U\)-linear map, written as
    \[
        1\longmapsto(f_0,\dots,f_n),
        \qquad f_0,\dots,f_n\in\mathcal{O}_X(U).
    \]
    Then \(\phi\) is a split injection if and only if there exist
    \(g_0,\dots,g_n\in\mathcal{O}_X(U)\) such that
    \[
        f_0g_0+\cdots+f_ng_n=1.
    \]
\end{proposition}

\begin{proof}
    \postdisplaypenalty=10000
    A map between finite free modules is determined by the images of the
    basis vectors. The same applies here to maps of sheaves of
    \(\mathcal{O}_U\)-modules. If \(e_0,\dots,e_n\) are the standard basis
    sections of \(\mathcal{O}_U^{\oplus(n+1)}\), then
    \[
        \phi(1)=(f_0,\dots,f_n)=\sum_{i=0}^n f_i e_i
    \]
    determines \(\phi\).

    Suppose first that \(\phi\) has a retraction
    \(\psi:\mathcal{O}_U^{\oplus(n+1)}\to\mathcal{O}_U\), so that
    \(\psi\circ\phi=\operatorname{id}_{\mathcal{O}_U}\).
    \[
    \begin{tikzcd}[column sep=large]
        \mathcal{O}_U \arrow[r,"\phi",bend left=15]
        & \mathcal{O}_U^{\oplus(n+1)}
            \arrow[l,"\psi",bend left=15]
    \end{tikzcd}
    \]
    Write \(\psi(e_i)=g_i\). Then
    \[
        1=\psi(\phi(1))
        =\sum_{i=0}^n f_i\psi(e_i)
        =\sum_{i=0}^n f_i g_i.
    \]

    Conversely, given such \(g_0,\dots,g_n\), define \(\psi\) on every open
    set \(V\subseteq U\) by
    \[
        \psi_V(a_0,\dots,a_n)
        =\sum_{i=0}^n (g_i|_V)a_i.
    \]
    These maps are \(\mathcal{O}_X(V)\)-linear and commute with restriction,
    so they define a morphism of sheaves of \(\mathcal{O}_U\)-modules.
    The identity \(\sum_i f_i g_i=1\) gives
    \(\psi\circ\phi=\operatorname{id}_{\mathcal{O}_U}\), as required.
\end{proof}

\begin{remark}\label{rem:lsi-sections}
    Here \(1\) denotes the unit section, and the \(f_i\) are sections of a
    sheaf, not fixed scalars. More precisely, for every open
    \(V\subseteq U\), the map on sections is
    \[
        \phi_V:\mathcal{O}_X(V)
        \longrightarrow\mathcal{O}_X(V)^{\oplus(n+1)},
        \qquad
        1_V\longmapsto(f_0|_V,\dots,f_n|_V).
    \]
    We will sometimes suppress these restriction symbols and simply write
    \(f_0,\dots,f_n\) when the open set is understood.
\end{remark}

\begin{corollary}\label{cor:lsi-unit-ideal}
    If \(\phi\) is a split injection as above, then
    \[
        (f_0,\dots,f_n)=\mathcal{O}_X(U).
    \]
    More generally, for every open \(V\subseteq U\),
    \[
        (f_0|_V,\dots,f_n|_V)=\mathcal{O}_X(V),
    \]
    and, for every \(x\in U\),
    \[
        (f_{0,x},\dots,f_{n,x})=\mathcal{O}_{X,x},
    \]
    where \(f_{i,x}\) denotes the germ of \(f_i\) at \(x\).
\end{corollary}

\begin{proof}
    Restrict the identity \(\sum_i f_i g_i=1\) to \(V\), or pass to its germ
    at \(x\).
\end{proof}

The stalk \(\mathcal{O}_{X,x}\) is a local ring, with unique maximal ideal
\(\mathfrak{m}_x\). Every principal ideal in this ring is either the whole
ring or is contained in \(\mathfrak{m}_x\). If all the \(f_{i,x}\) belonged
to \(\mathfrak{m}_x\), every linear combination
\(\sum_i g_{i,x}f_{i,x}\) would also belong to \(\mathfrak{m}_x\), and could
not equal 1. Consequently, at least one \(f_{i,x}\) lies outside
\(\mathfrak{m}_x\), and is therefore a unit.

\begin{corollary}\label{cor:lsi-unit-at-stalks}
    If \(\phi:\mathcal{O}_U\to\mathcal{O}_U^{\oplus(n+1)}\), with
    \(1\mapsto(f_0,\dots,f_n)\), is locally split, then for every \(x\in U\)
    the stalk map is
    \[
        \phi_x:\mathcal{O}_{X,x}
        \longrightarrow\mathcal{O}_{X,x}^{\oplus(n+1)},
        \qquad 1\longmapsto(f_{0,x},\dots,f_{n,x}),
    \]
    and at least one \(f_{i,x}\) is a unit in \(\mathcal{O}_{X,x}\).
\end{corollary}

\begin{proof}
    For each \(x\), choose a neighbourhood on which \(\phi\) splits, and
    apply the preceding argument there.
\end{proof}

A stalk consists of germs of sections near a point. If a germ \(f_{i,x}\)
is invertible, its inverse is represented by a section on some
neighbourhood of \(x\). After shrinking that neighbourhood, the product
of the two sections is 1. Thus \(f_i\) is invertible on some open
neighbourhood \(W\) of \(x\). This lets us recover information on a
neighbourhood from information at a stalk, and gives the following stalk
criterion.

\begin{proposition}\label{prop:lsi-stalk-criterion}
    Let \(\mathcal{L}\) be a locally free \(\mathcal{O}_X\)-module of rank 1.
    A map
    \(\iota:\mathcal{L}\to\mathcal{O}_X^{\oplus(n+1)}\) is locally split
    if and only if, for every \(x\in X\), the stalk map
    \[
        \iota_x:\mathcal{L}_x
        \longrightarrow\mathcal{O}_{X,x}^{\oplus(n+1)}
    \]
    is a split injection.
\end{proposition}

\begin{proof}
    If \(\iota\) splits on a neighbourhood of \(x\), taking stalks of the
    retraction gives a retraction of \(\iota_x\).

    Conversely, suppose every \(\iota_x\) is split. Fix \(x\) and choose
    a trivialisation of \(\mathcal{L}\) on a neighbourhood \(U\), writing
    \(\iota\) there as \(1\mapsto(f_0,\dots,f_n)\). A retraction of the
    stalk map gives coefficients \(g_i\in\mathcal{O}_{X,x}\) with
    \(\sum_i f_{i,x}g_i=1\). Since \(\mathcal{O}_{X,x}\) is local, at least
    one \(f_{i,x}\) is a unit. Shrink to an open neighbourhood
    \(W\subseteq U\) of \(x\) on which \(f_i\) is invertible. Then
    \[
        (a_0,\dots,a_n)\longmapsto (f_i|_W)^{-1}a_i
    \]
    defines a retraction in this trivialisation. These neighbourhoods
    \(W\), as \(x\) varies, form an open cover on which \(\iota\) splits.
\end{proof}

This still does not quite express our original expectation: at every
point, a locally split injection should give an inclusion of a
one-dimensional vector space into an \((n+1)\)-dimensional vector space,
without becoming zero. The following examples explain the distinction
between a stalk map and a fibre map.

\Needspace{24\baselineskip}
\begin{example}\label{ex:lsi-23}
    Consider the map of \(\Z\)-modules
    \[
        \phi:\Z\longrightarrow\Z^2,
        \qquad 1\longmapsto(2,3).
    \]
    The associated morphism of sheaves on \(\Spec\Z\), which we also
    denote by \(\phi\), is locally split. Indeed,
    \(D(2)\cup D(3)=\Spec\Z\), and we have the following retractions:
    \begin{itemize}
        \item On \(D(2)\), the element 2 is invertible, and a retraction is
        \[
            \Z[1/2]^2\longrightarrow\Z[1/2],
            \qquad (a,b)\longmapsto a/2.
        \]
        \item On \(D(3)\), the element 3 is invertible, and a retraction is
        \[
            \Z[1/3]^2\longrightarrow\Z[1/3],
            \qquad (a,b)\longmapsto b/3.
        \]
    \end{itemize}
    At the point \(x=(7)\), the fibre map is
    \[
        \overline{\phi}_x:\F_7\longrightarrow\F_7^2,
        \qquad 1\longmapsto(\overline{2},\overline{3})\ne(0,0).
    \]
    It is non-zero, so it embeds a one-dimensional vector space as a
    line in \(\F_7^2\).
\end{example}

\begin{example}\label{ex:lsi-46}
    Now consider
    \[
        \psi:\Z\longrightarrow\Z^2,
        \qquad 1\longmapsto(4,6).
    \]
    This is an injection, but it is not split, since
    \((4,6)=2\Z\ne\Z\). Its associated sheaf morphism is also injective,
    but is not locally split. To see what happens, compare the stalk
    map and the fibre map at \(x=(2)\).

    The stalk map is
    \[
        \psi_x:\Z_{(2)}\longrightarrow\Z_{(2)}^2,
        \qquad 1\longmapsto(4/1,6/1).
    \]
    Although \(4/1\) and \(6/1\) lie in the maximal ideal
    \(2\Z_{(2)}\), they have not been reduced modulo that ideal.
    The ring \(\Z_{(2)}\) is a domain, so \(4a=0\) implies \(a=0\).
    Hence \(\psi_x\) remains injective. Neither coordinate is a unit,
    however, so this stalk map is not split.

    The fibre map is instead
    \[
        \overline{\psi}_x:\F_2\longrightarrow\F_2^2,
        \qquad 1\longmapsto(\overline{4},\overline{6})=(0,0).
    \]
    Thus the fibre map is zero.
\end{example}

These examples show that injectivity of every stalk map is not enough
to imply that a sheaf morphism is locally split. The stalk criterion in
\Cref{prop:lsi-stalk-criterion} requires \emph{split} injections. Fibre
maps, on the other hand, give a criterion using injectivity alone.

\begin{definition}[Fibre map]\label{def:lsi-fibre-map}
    For \(x\in X\), let \(A=\mathcal{O}_{X,x}\) and let
    \(k(x)=A/\mathfrak{m}_x\) be the residue field at \(x\).
    A morphism \(\iota:\mathcal{L}\to\mathcal{O}_X^{\oplus(n+1)}\)
    induces an \(A\)-linear stalk map
    \(\iota_x:\mathcal{L}_x\to A^{\oplus(n+1)}\).
    Reducing modulo \(\mathfrak{m}_x\) gives the \emph{fibre map}
    \[
    \begin{tikzcd}[column sep=large]
        \mathcal{L}_x
            \arrow[r,"\iota_x"] \arrow[d,two heads]
        & A^{\oplus(n+1)} \arrow[d,two heads] \\
        \mathcal{L}_x/\mathfrak{m}_x\mathcal{L}_x
            \arrow[r,"\overline{\iota}_x"']
        & k(x)^{\oplus(n+1)}.
    \end{tikzcd}
    \]
    The bottom map is well-defined because \(\iota_x\) is \(A\)-linear,
    and therefore
    \[
        \iota_x(\mathfrak{m}_x\mathcal{L}_x)
        \subseteq\mathfrak{m}_x A^{\oplus(n+1)}.
    \]
\end{definition}

\begin{proposition}[Fibre criterion]\label{prop:lsi-fibre-criterion}
    Let \(\mathcal{L}\) be locally free of rank 1. A morphism
    \(\iota:\mathcal{L}\to\mathcal{O}_X^{\oplus(n+1)}\) is locally split
    if and only if, for every \(x\in X\), the fibre map
    \(\overline{\iota}_x\) is non-zero. Equivalently, every fibre map is
    injective.
\end{proposition}

\begin{proof}
    Suppose first that \(\iota\) is locally split. Near each \(x\), choose
    a retraction \(\rho\) with
    \(\rho\circ\iota=\operatorname{id}_{\mathcal{L}}\).
    Taking stalks and then reducing modulo \(\mathfrak{m}_x\) gives
    \[
        \overline{\rho}_x\circ\overline{\iota}_x
        =\operatorname{id}_{\mathcal{L}_x/
            \mathfrak{m}_x\mathcal{L}_x}.
    \]
    Hence \(\overline{\iota}_x\) is injective. Its source is a
    one-dimensional vector space, so it is also non-zero.

    Conversely, suppose \(\overline{\iota}_x\ne0\) for every \(x\).
    Choose a trivialisation on a neighbourhood \(U\) of \(x\), so that
    the map is written as
    \[
        \phi:\mathcal{O}_U\longrightarrow
        \mathcal{O}_U^{\oplus(n+1)},
        \qquad 1\longmapsto(f_0,\dots,f_n).
    \]
    Its fibre map, in this trivialisation, is
    \[
        \overline{\phi}_x:k(x)\longrightarrow k(x)^{\oplus(n+1)},
        \qquad 1\longmapsto\bigl(f_0(x),\dots,f_n(x)\bigr),
    \]
    where \(f_i(x)\) is the image of the germ \(f_{i,x}\) in \(k(x)\).
    Since this map is non-zero, some \(f_i(x)\ne0\). But
    \[
        f_i(x)\ne0
        \quad\Longleftrightarrow\quad
        f_{i,x}\notin\mathfrak{m}_x
        \quad\Longleftrightarrow\quad
        f_{i,x}\in\mathcal{O}_{X,x}^{\times}.
    \]
    This germ has an inverse, so \(f_i\) is invertible on some smaller
    neighbourhood \(W\) of \(x\). On \(W\), the map
    \((a_0,\dots,a_n)\mapsto(f_i|_W)^{-1}a_i\) is a retraction in the
    chosen trivialisation. Such neighbourhoods give an open cover on
    which \(\iota\) splits.

    Finally, a linear map from a one-dimensional vector space is
    injective if and only if it is non-zero.
\end{proof}

\begin{remark}\label{rem:lsi-fibre-line}
    The residue ring \(k(x)\) is a field. The quotient
    \(\mathcal{L}_x/\mathfrak{m}_x\mathcal{L}_x\), however, is naturally
    a one-dimensional \(k(x)\)-vector space, not a field with a chosen
    element 1. A trivialisation gives an isomorphism
    \[
        \mathcal{L}_x/\mathfrak{m}_x\mathcal{L}_x\cong k(x),
    \]
    but this isomorphism depends on the choice of trivialisation.
    The non-zero fibre map embeds this one-dimensional vector space
    into \(k(x)^{n+1}\). This is exactly the line at \(x\) that we wanted,
    in agreement with the description of maps
    \(\Spec k\to\P^n_\Z\).
\end{remark}

\Needspace{22\baselineskip}
We can now collect the equivalent criteria.

\begin{proposition}\label{prop:lsi-equivalent-conditions}
    Let \(\mathcal{L}\) be a locally free \(\mathcal{O}_X\)-module of rank 1
    and let \(\iota:\mathcal{L}\to\mathcal{O}_X^{\oplus(n+1)}\) be a
    morphism. The following are equivalent:
    \begin{enumerate}
        \item \(\iota\) is locally split.
        \item For every \(x\in X\), the stalk map \(\iota_x\) is a split
        injection.
        \item For every \(x\in X\), after choosing a local trivialisation
        near \(x\) and writing \(1\mapsto(f_0,\dots,f_n)\), at least one
        germ \(f_{i,x}\) is a unit in \(\mathcal{O}_{X,x}\).
        \item For every \(x\in X\), the fibre map
        \(\overline{\iota}_x\) is non-zero.
        \item For every \(x\in X\), the fibre map
        \(\overline{\iota}_x\) is injective.
    \end{enumerate}
\end{proposition}

\begin{proof}
    The stalk criterion and its proof give the equivalence of the first
    three conditions. The fibre criterion identifies the first condition
    with the fourth and fifth.
\end{proof}

\subsection*{Constructing a morphism to projective space}

We now use the preceding criteria to construct a morphism to projective
space from a line bundle with a locally split injection.

\begin{proposition}\label{prop:lsi-map-to-projective-space}
    Let \(\mathcal{L}\) be a line bundle on \(X\). A locally split injection
    \[
        \iota:\mathcal{L}\hookrightarrow
        \mathcal{O}_X^{\oplus(n+1)}
    \]
    canonically determines a morphism
    \(F:X\to\P^n_\Z\).
\end{proposition}

\begin{proof}
    We construct the morphism locally and check agreement on overlaps.
    Write
    \[
        V_i=\Spec\Z\left[\frac{X_r}{X_i}\;\middle|\;r\ne i\right]
        \subseteq\P^n_\Z
    \]
    for the standard affine charts.

    \emph{Local construction.}
    Choose \(x\in X\), a neighbourhood \(U\), and a trivialisation
    \(\alpha:\mathcal{O}_U\xrightarrow{\sim}\mathcal{L}|_U\).
    In these coordinates, write
    \[
        \phi=(\iota|_U)\circ\alpha:
        \mathcal{O}_U\longrightarrow\mathcal{O}_U^{\oplus(n+1)},
        \qquad 1\longmapsto(f_0,\dots,f_n).
    \]
    At least one \(f_{i,x}\) is a unit. Suppose it is \(f_{0,x}\), and
    shrink \(U\) so that \(f_0\in\mathcal{O}_X(U)^\times\).
    The ring homomorphism
    \[
        \Z\left[\frac{X_1}{X_0},\dots,\frac{X_n}{X_0}\right]
        \longrightarrow\mathcal{O}_X(U),
        \qquad \frac{X_r}{X_0}\longmapsto\frac{f_r}{f_0}
    \]
    defines \(q_0:U\to V_0\), hence a morphism
    \[
        F_U:U\xrightarrow{q_0}V_0\hookrightarrow\P^n_\Z.
    \]
    Here \(U\) need not be affine: we use that the target \(V_0\) is affine.

    \emph{Independence of choices.}
    First keep \(U\) fixed and change the trivialisation. Any other
    trivialisation is \(\alpha'=\alpha\circ(\times g)\) for some
    \(g\in\mathcal{O}_X(U)^\times\). This replaces the coordinates by
    \((gf_0,\dots,gf_n)\), but
    \[
        \frac{gf_r}{gf_0}=\frac{f_r}{f_0},
    \]
    so the ring map, and therefore \(F_U\), is unchanged.

    There is also the choice of unit coordinate. Suppose \(f_0\) and
    \(f_1\) are both units on \(U\). Choosing \(f_1\) instead gives
    \(q_1:U\to V_1\), corresponding to
    \(X_r/X_1\mapsto f_r/f_1\). We must compare their composites to
    \(\P^n_\Z\). Put
    \[
        V_{01}=D\left(\frac{X_1}{X_0}\right)\subseteq V_0,
        \qquad
        V_{10}=D\left(\frac{X_0}{X_1}\right)\subseteq V_1.
    \]
    These opens are identified by the transition isomorphism
    \(\tau_{01}:V_{01}\xrightarrow{\sim}V_{10}\) used to construct
    \(\P^n_\Z\). Since \(f_1/f_0\) and \(f_0/f_1\) are units, the two
    maps factor through these opens:
    \[
    \begin{tikzcd}[column sep=large,row sep=large]
        & V_{01}
            \arrow[r,hook]
            \arrow[dd,"\tau_{01}"',"\sim"]
        & V_0 \arrow[dr,hook] & \\
        U \arrow[ur,"\widetilde q_0"]
          \arrow[dr,"\widetilde q_1"']
        & & & \P^n_\Z \\
        & V_{10} \arrow[r,hook]
        & V_1 \arrow[ur,hook] &
    \end{tikzcd}
    \]
    To check commutativity, compare the ring maps. For example, when
    \(n\ge2\), following the transition and then the map from \(U\) gives
    \[
        \frac{X_2}{X_1}
        \longmapsto\frac{X_2}{X_0}\frac{X_0}{X_1}
        \longmapsto\frac{f_2}{f_0}\frac{f_0}{f_1}
        =\frac{f_2}{f_1}.
    \]
    This is its image under the ring map for \(\widetilde q_1\).
    The other generators are checked in the same way, so
    \(\widetilde q_1=\tau_{01}\circ\widetilde q_0\).
    Thus the two composites to \(\P^n_\Z\) agree. The same argument
    applies to any two unit coordinates.

    For \(\P^1\), this is just the choice between slope \(f_1/f_0\)
    and anti-slope \(f_0/f_1\). They are reciprocals, exactly as in
    the gluing of the two charts of \(\P^1\).

    \emph{Gluing on \(X\).}
    These neighbourhoods cover \(X\). On \(U_1\cap U_2\), the restrictions
    of the local morphisms are obtained by restricting their
    trivialisations and unit coordinates. These give two choices for
    the same data \((\mathcal{L},\iota)|_{U_1\cap U_2}\). Since neither
    choice affects the resulting morphism,
    \[
        F_{U_1}|_{U_1\cap U_2}=F_{U_2}|_{U_1\cap U_2}.
    \]
    The local morphisms therefore glue uniquely to a morphism
    \(F:X\to\P^n_\Z\). The same overlap argument shows that it is
    independent of the chosen neighbourhoods.
\end{proof}

% Third part

\subsection*{The reverse construction}

Conversely, suppose we are given a morphism
\(\varphi:X\to\P^n_\Z\). Let
\[
    V_i=\Spec\Z\left[
        \frac{X_r}{X_i}\;\middle|\;r\ne i
    \right],
    \qquad
    U_i=\varphi^{-1}(V_i).
\]
The \(U_i\) form an open cover of \(X\). For each \(i\), the
restriction \(\varphi|_{U_i}:U_i\to V_i\) gives a ring map
\[
    \Z\left[
        \frac{X_r}{X_i}\;\middle|\;r\ne i
    \right]
    \longrightarrow\mathcal O_X(U_i),
    \qquad
    \frac{X_r}{X_i}\longmapsto h_r^{(i)}.
\]
Set \(h_i^{(i)}=1\). On \(U_i\), take the trivial line bundle
and define
\[
    \iota_i:\mathcal O_{U_i}
    \longrightarrow\mathcal O_{U_i}^{\oplus(n+1)},
    \qquad
    1\longmapsto
    \bigl(h_0^{(i)},\ldots,h_n^{(i)}\bigr).
\]
Projection onto the \(i\)-th coordinate is a left inverse, so
\(\iota_i\) is a split injection. We now need to glue these
local line bundles and maps.

Consider first \(U_{01}=U_0\cap U_1\). Write the restrictions
of \(\iota_0\) and \(\iota_1\) as
\[
    \begin{aligned}
        h:\mathcal O_{U_{01}}
        &\longrightarrow\mathcal O_{U_{01}}^{\oplus(n+1)},
        &1&\longmapsto(1,h_1,h_2,\ldots,h_n),\\
        h':\mathcal O_{U_{01}}
        &\longrightarrow\mathcal O_{U_{01}}^{\oplus(n+1)},
        &1&\longmapsto(h'_0,1,h'_2,\ldots,h'_n).
    \end{aligned}
\]
Here \(h_r\) is the pullback of \(X_r/X_0\), while \(h'_r\)
is the pullback of \(X_r/X_1\), with all sections restricted
to \(U_{01}\). These vectors need not be equal. However,
the coordinate relations on projective space give
\[
    h'_0=h_1^{-1}.
\]
For example, when \(n\ge2\),
\[
    \frac{X_2}{X_1}
    =
    \frac{X_2}{X_0}\frac{X_0}{X_1}
    \quad\Longrightarrow\quad
    h'_2=h_2h_1^{-1}.
\]
The other coordinates behave in the same way. Hence
\[
    h'(1)
    =
    (h_1^{-1},1,h_2h_1^{-1},\ldots,h_nh_1^{-1})
    =
    h_1^{-1}h(1).
\]
Thus we should identify the two sources by multiplication
by \(h_1\), rather than by the identity. This makes the
following diagram commute:
\[
\begin{tikzcd}[column sep=large,row sep=large]
    \mathcal O_{U_{01}}
        \arrow[r,"h"]
        \arrow[d,"\times h_1"']
    & \mathcal O_{U_{01}}^{\oplus(n+1)}
        \arrow[d,equal]\\
    \mathcal O_{U_{01}}
        \arrow[r,"h'"']
    & \mathcal O_{U_{01}}^{\oplus(n+1)}.
\end{tikzcd}
\]
Indeed, \(h'\circ(\times h_1)=h\), and multiplication by
\(h_1\) is an isomorphism since \(h_1\) is a unit on
\(U_{01}\).

More generally, on \(U_{ij}=U_i\cap U_j\), define
\[
    h_{ij}
    =
    \varphi^\#\left(\frac{X_j}{X_i}\right)
    \in\mathcal O_X(U_{ij})^\times.
\]
Its inverse is \(h_{ji}\). The same calculation gives
\[
    \iota_j\circ(\times h_{ij})=\iota_i
    \qquad\text{on }U_{ij}.
\]
On a triple overlap, the coordinate ratios satisfy
\[
    \frac{X_j}{X_i}
    \frac{X_k}{X_j}
    \frac{X_i}{X_k}=1,
\]
so their pullbacks satisfy the cocycle condition
\[
    h_{ij}h_{jk}h_{ki}=1.
\]
We may therefore glue the sheaves \(\mathcal O_{U_i}\)
using the isomorphisms
\[
    \times h_{ij}:
    \mathcal O_{U_i}|_{U_{ij}}
    \xrightarrow{\sim}
    \mathcal O_{U_j}|_{U_{ij}}.
\]
This gives a line bundle \(\mathcal L\), since
\(\mathcal L|_{U_i}\cong\mathcal O_{U_i}\).
The compatibility of the \(\iota_i\) with these
identifications makes them glue to a global map
\[
    \iota:\mathcal L\longrightarrow
    \mathcal O_X^{\oplus(n+1)}.
\]
It is split on each \(U_i\), and hence is a locally split
injection.

If
\(\alpha_i:\mathcal O_{U_i}\xrightarrow{\sim}\mathcal L|_{U_i}\)
are the resulting trivialisations, the gluing relation is
\[
    \alpha_i=\alpha_j\circ(\times h_{ij}).
\]
Thus, in the convention \(\alpha_j=\lambda_{ij}\alpha_i\),
we have \(\lambda_{ij}=h_{ij}^{-1}\).

Locally, the two constructions undo one another: the
normalised vectors recover the affine coordinate ratios,
and the ratios recover the same local line and injection.
We therefore obtain the correspondence
\[
    \operatorname{Mor}(X,\P^n_\Z)
    \longleftrightarrow
    \{(\mathcal L,\iota)\}/\cong.
\]

\begin{remark}
    The local sources are trivial, but their identifications
    on overlaps need not be the identity. The resulting
    \(\mathcal L\) need not be globally trivial: a family of
    lines may have local generating vectors without a global
    generating vector. This is why we need a line bundle
    rather than always taking \(\mathcal L=\mathcal O_X\).

    The same construction works for a locally ringed space
    \(X\).
\end{remark}

\end{document}